% Generated draft: requires mathematical and source review.
\documentclass[11pt]{article}
\usepackage[margin=25mm]{geometry}
\usepackage{fontspec}
\setmainfont{Latin Modern Roman}
\usepackage{amsmath,amssymb,mathtools,unicode-math}
\setmathfont{Latin Modern Math}
\usepackage{enumitem,needspace,xurl,hyperref}
\setlength{\parindent}{0pt}
\setlength{\parskip}{5pt}
\setlength{\emergencystretch}{3em}
\newcommand{\N}{\mathbb{N}}
\newcommand{\Z}{\mathbb{Z}}
\newcommand{\Q}{\mathbb{Q}}
\newcommand{\R}{\mathbb{R}}
\newcommand{\C}{\mathbb{C}}
\newcommand{\F}{\mathbb{F}}
\newcommand{\A}{\mathbb{A}}
\newcommand{\PP}{\mathbb{P}}
\newcommand{\E}{\mathbb{E}}
\newcommand{\eps}{\varepsilon}
\newcommand{\dd}{\,\mathrm{d}}
\begin{document}
\section*{UnsolvedMath ID: 30000565}
\textbf{Source problem code:} catalogue entry
\subsection{Definitions and mathematical statement}
A concordance between two ordered links in $S^3$ is a collection of disjoint annuli in $S^3\times[0,1]$ joining corresponding components, in the chosen smooth or locally flat category. A Whitney disk pairs a pair of transverse intersection points of opposite sign by a disk whose boundary consists of arcs on the two intersecting sheets; its interior may be immersed and meet other sheets. A Whitney tower recursively provides such pairings, with the relevant framing and intersection-tree conventions.

The general question is whether the proposed higher-order Arf constructions from Whitney towers give invariants independent of tower choices and constant under link concordance. Their precise target groups and correction or indeterminacy relations are not given in the extract.

The specified first test case is the two-component Bing double $L=l_1\cup l_2$ of the figure-eight knot. Retain the source's disks
\[
\partial D_i=l_i,\qquad D_i\subset B^4\quad(i=1,2),
\]
each embedded individually; they are not assumed mutually disjoint. An immersed Whitney disk $W_{(1,2)}$ pairs intersections between $D_1$ and $D_2$. A second Whitney disk $W_{((1,2),(1,2))}$ pairs self-intersections of $W_{(1,2)}$. The indicated configuration has a single transverse interior intersection between these two Whitney disks. The question is whether that intersection represents a well-defined, nontrivial concordance obstruction for $L$, rather than an artifact of the chosen configuration. Nontriviality would require a defined invariant that detects $L$ nontrivially, not merely the presence of an intersection in one tower.

With order-zero initial disks and the convention that a Whitney disk pairing sheets of orders $a,b$ has order $a+b+1$, the two Whitney disks have orders $1$ and $3$, while their intersection has order $4$. The original extract calls the indicated intersection ``order-3''; this may refer to the higher Whitney disk, but the convention is not resolved here. Neither raw intersection parity nor a particular quotient of tree-valued intersections is silently declared to be the desired invariant.
\subsection{Short English statement}
Ask whether Whitney-tower higher Arf constructions are concordance invariants, and whether the specified single-intersection configuration for the Bing double of the figure-eight knot detects a nontrivial invariant. Tower choices, framing relations and the order convention require clarification.
\subsection{Sources}
\begin{itemize}
\item Supplied report: \emph{Four-dimensional Manifolds} (2007), Oberwolfach Reports, DOI 10.4171/owr/2006/35, pp. 2141--2194; proposed year 2006. \url{https://doi.org/10.4171/owr/2006/35}.
\item Local UnsolvedMath registry, permanent ID 30000565, catalogue entry; created 25 August 2026, updated 16 September 2026. Its corrected\_verified note says the report identifies the Bing double, both bounding disks, both Whitney disks and the single interior intersection. The supplied 21 August 2026 literature triage labels the extracted claim open and reports no located authoritative resolution. No additional external reference is supplied; the dated assessment has not been independently verified here.
\end{itemize}
\end{document}
